\documentclass[11pt]{article}
\usepackage[utf8]{inputenc}
\usepackage[T1]{fontenc}
\usepackage{lmodern}
\usepackage{amsmath,amssymb,amsthm}
\usepackage[margin=1in]{geometry}
\usepackage{booktabs,tabularx}
\usepackage[hidelinks]{hyperref}
\usepackage{url}
\usepackage{array}
\def\UrlBreaks{\do0\do1\do2\do3\do4\do5\do6\do7\do8\do9\do a\do b\do c\do d\do e\do f\do/\do.\do-}
\emergencystretch=3em
\newtheorem{theorem}{Theorem}
\newtheorem{lemma}[theorem]{Lemma}
\newtheorem{corollary}[theorem]{Corollary}
\newtheorem{conjecture}[theorem]{Conjecture}
\theoremstyle{remark}\newtheorem{remark}[theorem]{Remark}
\newcommand{\lean}[1]{{\ttfamily\def\_{\textunderscore\allowbreak}\def\/{/\allowbreak}#1}}
\author{Hetark Shah, AI-assisted by Grok}
\newcommand{\Hh}{H^{(2)}}
\date{October 12, 2026}
\title{Wolstenholme numbers and a gcd conjecture of Eldar and Ordowski (OEIS A007406): infinite families, formalized in Lean~4}
\begin{document}
\maketitle
\begin{abstract}
Let $a(n)$ be the numerator of $\Hh_n=\sum_{k=1}^n k^{-2}$ (OEIS A007406~\cite{oeis}). Eldar and Ordowski conjectured in 2019 that $\gcd(n,a(n-1))$ equals $n$ for primes $n>3$ and equals $1$ for composites $n>3$. \textbf{The full conjecture remains open.} We prove a leading-digit criterion: if $q\le m$ is prime and $q$ does not divide $a(U)$, where $U$ is the leading base-$q$ digit of $m$, then $q\nmid a(m)$. As consequences we show that $\gcd(6,a(m))=1$ for all $m\ge1$, and that the conjecture holds for every $n>3$ that is prime (via Wolstenholme's theorem, which we formalize), $3$-smooth, of the form $2q^k$ with $q$ an odd prime, and more generally $cq^k$ with $c\ge2$ $3$-smooth and $q>c$ a prime not dividing $a(c-1)$. All results are formalized in Lean~4 with Mathlib against the definitions in \texttt{formal-conjectures}~\cite{fc}, with no \texttt{sorry} and only the standard axioms.
\end{abstract}

\section{Introduction}
OEIS A007406~\cite{oeis} contains the comment: \emph{``Conjecture: for $n>3$, $\gcd(n,a(n-1))=\mathrm{A089026}(n)$. Checked up to $n=10^5$.''} (A.~Eldar and T.~Ordowski, Jul.~28, 2019). Here $\mathrm{A089026}(n)=n$ if $n$ is prime and $1$ otherwise. J.~Sondow added that the prime case follows from Wolstenholme's theorem and that \emph{``it remains to show that $\gcd(n,a(n-1))=1$ if $n>3$ is composite.''} The conjecture is formalized as \lean{OeisA7406.conjecture} in \nolinkurl{FormalConjectures/OEIS/7406.lean} of \texttt{formal-conjectures}~\cite{fc}, where it is tagged \texttt{research open}.

\section{Related work}
Our key lemma is not new in substance. It is the order-$2$ analogue of the classical leading-digit argument for the sets $J(p)=\{n:p\mid\text{numerator of }H_n\}$, studied by Eswarathasan and Levine~\cite{el}, Boyd~\cite{boyd}, Sanna~\cite{sanna} and Wu and Chen~\cite{wc}; Kamano~\cite{kam} treats $3$-adic valuations of generalized harmonic numbers. Wolstenholme's theorem and its refinements are surveyed by Me\v{s}trovi\'c~\cite{mes}. As of October 12, 2026, the OEIS still lists the conjecture as open, and we found no work applying these tools to it. Our contribution is the reduction to a leading-digit criterion, the explicit infinite families, and machine-checked proofs stated against the upstream definitions (\texttt{formal-conjectures} commit \texttt{34ee0e3}). The formalization uses Lean~4~\cite{lean4} and Mathlib~\cite{mathlib}.

\section{Statement}
\begin{theorem}\label{thm:main}
The equality $\gcd(n,a(n-1))=\mathrm{A089026}(n)$ holds for every $n>3$ that satisfies one of the following:
\begin{enumerate}
\item $n$ is prime;
\item $n=2^i3^j$;
\item $n=cq^k$ with $c=2^i3^j\ge2$, $k\ge1$, and $q>c$ a prime with $q\nmid a(c-1)$. This includes $n=2q^k$ for every odd prime $q$, $n=3q$ for $q\ge7$, and $n=4q$ for $q\ge5$, $q\ne7$.
\end{enumerate}
\end{theorem}
\noindent Cases not covered include $n=q^k$ with $k\ge2$, and general $n=pq$ in which neither cofactor is $3$-smooth.

\section{Proof overview}
\subsection{The leading-digit lemma}
Throughout, $q$ is prime and $v_q$ is the $q$-adic valuation on $\mathbb{Q}$.
\begin{lemma}\label{lem:key}
Let $m\ge1$, $s=\lfloor\log_q m\rfloor$ and $U=\lfloor m/q^s\rfloor$, so $1\le U<q$. If $q\nmid a(U)$ then $v_q(\Hh_m)=-2s$.
\end{lemma}
\begin{proof}
Since $U<q$, every term $u^{-2}$, $1\le u\le U$, is a $q$-adic unit, so $v_q(\Hh_U)\ge0$; as $q\nmid a(U)$ and the denominator of $\Hh_U$ is coprime to $q$, $v_q(\Hh_U)=0$. Split $\{1,\dots,m\}$ into $S_1=\{k:q^s\mid k\}=\{q^su:1\le u\le U\}$ and $S_0$. Then $\sum_{S_1}k^{-2}=q^{-2s}\Hh_U$ has valuation $-2s$. For $k\in S_0$ we have $v_q(k)\le s-1$, so $v_q(k^{-2})\ge -2s+2$; since all terms are positive, the ultrametric inequality gives $v_q(\sum_{S_0}k^{-2})\ge-2s+1$. Hence $v_q(\Hh_m)=-2s$.
\end{proof}
\begin{theorem}[criterion]\label{thm:crit}
If $q\le m$ and $q\nmid a(\lfloor m/q^{\lfloor\log_q m\rfloor}\rfloor)$, then $q\nmid a(m)$. Consequently, if $n>3$ is composite and for every prime $q\mid n$, $q$ does not divide $a$ of the leading base-$q$ digit of $n-1$, then $\gcd(n,a(n-1))=1$.
\end{theorem}
\begin{proof}
Here $s\ge1$, so $v_q(\Hh_m)<0$; if $q\mid a(m)$ then $v_q(\Hh_m)\ge1$. For the second claim, a prime divisor $q$ of composite $n$ satisfies $q<n$, so $q\le n-1$.
\end{proof}
\subsection{Consequences}
\begin{theorem}\label{thm:six}
For all $m\ge1$, $\gcd(6,a(m))=1$.
\end{theorem}
\begin{proof}
For $q=2$ the leading digit is $1$ and $a(1)=1$; for $q=3$ it is $1$ or $2$, and $a(1)=1$, $a(2)=5$. Small $m<q$ are checked directly.
\end{proof}
\begin{corollary}
The conjecture holds for every $n=2^i3^j>3$.
\end{corollary}
\begin{theorem}\label{thm:cq}
Let $c=2^i3^j\ge2$ and let $q>c$ be prime with $q\nmid a(c-1)$. Then $\gcd(cq,a(cq-1))=1$, that is, the conjecture holds at $n=cq$.
\end{theorem}
\begin{proof}
Primes $2,3$ are handled by Theorem~\ref{thm:six}. For $q$: $q\le cq-1<q^2$, so $s=1$ and the leading digit is $\lfloor (cq-1)/q\rfloor=c-1$.
\end{proof}
\begin{corollary}
The conjecture holds for $n=2q$ (all primes $q\ge3$), $n=3q$ (primes $q\ge7$), and $n=4q$ (primes $q\ge5$, $q\neq7$), using $a(1)=1$, $a(2)=5$, $a(3)=49$. For each $3$-smooth $c\ge2$ it holds for $n=cq$ for all primes $q>\max(c,a(c-1))$.
\end{corollary}

\begin{theorem}\label{thm:cqk}
Let $c=2^i3^j\ge2$, $q>c$ prime with $q\nmid a(c-1)$, and $k\ge1$. Then the conjecture holds at $n=cq^k$. In particular it holds for $n=2q^k$, for every odd prime $q$ and every $k\ge1$.
\end{theorem}
\begin{proof}
For $m=cq^k-1$ we have $q^k\le(c-1)q^k\le m<cq^k<q^{k+1}$, so $\lfloor\log_q m\rfloor=k$ and the leading digit is $c-1$. Primes $2,3$ are handled by Theorem~\ref{thm:six}.
\end{proof}

\subsection{The prime direction}
\begin{theorem}[Wolstenholme]\label{thm:w}
For every prime $p\ge5$, $p\mid a(p-1)$. Hence $\gcd(p,a(p-1))=p=\mathrm{A089026}(p)$ for every prime $p>3$.
\end{theorem}
\begin{proof}
Let $F=(p-1)!$ and $N=\sum_{k=1}^{p-1}(F/k)^2$, so that $\Hh_{p-1}F^2=N$. In $\mathbb{F}_p$, Wilson gives $k\cdot(F/k)=-1$, so $(F/k)^2=(F/k)^2k^{p-1}=k^{p-3}$. Therefore $N\equiv\sum_{x\in\mathbb{F}_p}x^{p-3}=0$, since $0<p-3<p-1$. As $p\nmid F^2$, $p$ divides the numerator.
\end{proof}
\begin{corollary}
The full conjecture $\gcd(n,a(n-1))=\mathrm{A089026}(n)$ holds for every $n>3$ that is prime, $3$-smooth, or of the form $2q^k$ ($q$ an odd prime, $k\ge1$), and for $n=cq^k$ as in Theorem~\ref{thm:cqk}.
\end{corollary}

\begin{remark}[Coverage and limits]\label{rem:cov}
A short computation (\texttt{coverage.py}) shows the criterion of Theorem~\ref{thm:crit} settles $76{,}100$ of the $90{,}407$ composites $4\le n\le10^5$ (about $84\%$). It cannot settle every case: $q\mid a(q-1)$ (Wolstenholme) and $q\mid a((q-1)/2)$ for all $q\ge5$, so for instance $n=q^2$ needs a second-order analysis. There $\Hh_{q^2-1}=\sum_{q\nmid k<q^2}k^{-2}+q^{-2}\Hh_{q-1}$, and the first sum is $\equiv0\pmod{q^2}$. One finds that $q\mid a(q^2-1)$ would force $v_q(\Hh_{q-1})\ge3$, which by Glaisher's congruence $\Hh_{q-1}\equiv\tfrac{2}{3}qB_{q-3}\pmod{q^2}$ requires $q\mid B_{q-3}$ (a Wolstenholme prime, e.g.\ $16843$) together with a further congruence. We do not formalize this.
\end{remark}


\section{Formalization}
All declarations are in namespace \lean{Wolst}. The definitions \lean{a} and \lean{A89026} in \lean{Basic.lean} are copied verbatim from \texttt{formal-conjectures}, and every main theorem concludes exactly \lean{Nat.gcd n (a (n - 1)) = A89026 n}.
\begin{table}[ht]
\centering
\begin{tabularx}{\linewidth}{@{}>{\raggedright\arraybackslash}X>{\raggedright\arraybackslash}X@{}}
\toprule
Lean name (file) & Statement \\
\midrule
\lean{padicValRat\_H}, \lean{not\_dvd\_a} (\lean{Basic.lean}) & Lemma~\ref{lem:key}; first part of Theorem~\ref{thm:crit} \\
\lean{conjecture\_of\_leading\_digits} (\lean{Families.lean}) & Theorem~\ref{thm:crit} \\
\lean{coprime\_six\_a}, \lean{conjecture\_smooth} & Theorem~\ref{thm:six} and its corollary \\
\lean{conjecture\_smooth\_mul\_prime} & Theorem~\ref{thm:cq} \\
\lean{conjecture\_two/three/four\_mul\_prime}, \lean{conjecture\_smooth\_mul\_large\_prime} & corollaries of Theorem~\ref{thm:cq} \\
\lean{conjecture\_smooth\_mul\_prime\_pow}, \lean{conjecture\_two\_mul\_prime\_pow} (\lean{Prime.lean}) & Theorem~\ref{thm:cqk} \\
\lean{prime\_dvd\_a}, \lean{conjecture\_prime} & Theorem~\ref{thm:w} \\
\lean{conjecture\_on\_families} & cases (1)--(2) and $2q^k$ of Theorem~\ref{thm:main} \\
\bottomrule
\end{tabularx}
\caption{Theorems and their Lean names.}\label{tab:lean}
\end{table}

\section{Reproducibility}
\begin{itemize}
\item Source code: \url{https://github.com/het4rk/eldar-ordowski-a007406}. The Lean project is in \texttt{Wolst/}.
\item Toolchain: \texttt{leanprover/lean4:v4.35.0-rc4}. Mathlib: tag \texttt{v4.35.0-rc4}, commit \nolinkurl{1f414401f69059aa7eead47b53ee40bd38455eeb}, pinned in \texttt{lake-manifest.json}.
\item Build:
\begin{verbatim}
cd Wolst
lake exe cache get
lake build
lake env lean Axioms.lean
\end{verbatim}

\item Axioms: each of the 15 declarations in \texttt{Axioms.lean} uses only the axioms \texttt{propext}, \texttt{Classical.choice} and \texttt{Quot.sound}.
\item Attestation: \texttt{verification/} contains the fresh-clone build log, the axiom output and \texttt{VERIFICATION.md}, which records the commit, the toolchain and SHA-256 hashes of the Lean sources. CI (\url{https://github.com/het4rk/eldar-ordowski-a007406/actions/workflows/lean.yml}) rebuilds the project on every push. It rejects \texttt{sorry} and non-whitelisted axioms, and uploads the build log and axiom report as artifacts.
\end{itemize}

\section{Limitations}
\textbf{The full conjecture remains open.} The criterion can never settle $n=q^2$, because $q\mid a(q-1)$; a second-order analysis would be needed (Remark~\ref{rem:cov}), and it is not formalized. The coverage figure in that remark was computed by \texttt{paper/coverage.py}. It is not a formal result.

\begin{thebibliography}{99}
\bibitem{oeis} OEIS Foundation Inc., \emph{A007406: Wolstenholme numbers}, \url{https://oeis.org/A007406}.
\bibitem{fc} Google DeepMind, \emph{formal-conjectures}, \url{https://github.com/google-deepmind/formal-conjectures}, file \nolinkurl{FormalConjectures/OEIS/7406.lean}.
\bibitem{el} A.~Eswarathasan, E.~Levine, $p$-integral harmonic sums, \emph{Discrete Math.} \textbf{91} (1991), 249--257.
\bibitem{boyd} D.~W.~Boyd, A $p$-adic study of the partial sums of the harmonic series, \emph{Experiment.\ Math.} \textbf{3} (1994), 287--302.
\bibitem{kam} K.~Kamano, On $3$-adic valuations of generalized harmonic numbers, \emph{Integers} \textbf{11} (2011).
\bibitem{sanna} C.~Sanna, On the $p$-adic valuation of harmonic numbers, \emph{J.\ Number Theory} \textbf{166} (2016), 41--46.
\bibitem{wc} B.-L.~Wu, Y.-G.~Chen, On certain properties of harmonic numbers, \emph{J.\ Number Theory} \textbf{175} (2017), 66--86.
\bibitem{mes} R.~Me\v{s}trovi\'c, Wolstenholme's theorem: its generalizations and extensions in the last hundred and fifty years, arXiv:1111.3057.
\bibitem{mathlib} The mathlib Community, \emph{The Lean mathematical library}, in Proc.\ CPP 2020, ACM, 367--381.
\bibitem{lean4} L.~de Moura, S.~Ullrich, \emph{The Lean 4 theorem prover and programming language}, in CADE-28, LNCS 12699, Springer, 2021, 625--635.
\end{thebibliography}
\end{document}
